\documentclass{article}

\usepackage{comment}
\usepackage{tabularx}
\usepackage{booktabs}
\usepackage{comment}

\setlength{\emergencystretch}{1em}

\title{Problem Statement and Goals\\\progname}

\author{\authname}

\date{}

\input{../Comments.text}
\input{../Common.text}

\begin{document}

\maketitle

\begin{table}[hp]
\caption{Revision History} \label{TblRevisionHistory}
\begin{tabularx}{\textwidth}{llX}
\toprule
\textbf{Date} & \textbf{Developer(s)} & \textbf{Change}\\
\midrule
September 26, 2026 & Shivya, Anindita, Harmanpreet, Vanessa& Initial version\\
% Date2 & Name(s) & Description of changes\\
... & ... & ...\\
\bottomrule
\end{tabularx}
\end{table}

\wss{You should check your problem statement with the
\href{https://smiths.github.io/capTemplate/Checklists/ProbState-Checklist.pdf}
{problem statement checklist}}.

\section{Problem Statement}

\begin{comment}
\wss{You should check your problem statement with the
\href{https://smiths.github.io/capTemplate/Checklists/ProbState-Checklist.pdf}
{problem statement checklist}}. 

\wss{You can change the section headings, as long as you include the required
information.}
\end{comment}

\textit{The Department of Computing and Software (CAS) lacks a centralized, coordinated
method for managing course outline schedules and ensuring academic policy/engineering accreditation compliance, 
resulting in congested workload periods for students and high administrative overhead for staff and faculty.}

\subsection{Problem}

\begin{comment}
    \wss{What problem are you trying to solve?  Why is it important? Do not talk
    about how you will solve the problem, only what the problem is.  AI is rarely
    part of the problem.  The problem is ``what'' you want to do, not ``how'' you
    want to do it.  The problem statement should be written in a way that is
    understandable to a non-technical audience.}
\end{comment}

Across the Department of Computing and Software (CAS), course outlines and assessment
schedules are currently created and managed individually by instructors.
As there is no centralized, coordinated system of determining when assignments and midterms
are scheduled, students in the same program and academic year frequently face severe workload
clustering. Having multiple midterms or major deliverables due within the same week or
on consecutive days results in increased academic stress and poorer learning outcomes for students.

\bigskip
From an administrative perspective, coordinating all of these schedules and ensuring institutional
guidelines are met is a time-consuming process. The CAS administrative
staff must manually inspect course outlines to verify compliance with departmental policies, university
regulations, and Canadian Engineering Accreditation Board (CEAB) criteria. This process is further
complicated by circumstances where missing or incomplete information is found in course outlines
(e.g., midterms listed as ``TBD''), resulting in administrative staff having to engage in back-and-forth
communication with instructors to clarify details. Furthermore, as course instructors have full academic
freedom over their course delivery, administrators are not authorized to reassign assignment deadlines or exam dates.

\bigskip
In summary, without a mechanism to identify scheduling bottlenecks and policy issues early, 
administrative workload remains high, staff and instructors are forced to spend time engaging in redundant communication,
and students continue to suffer from congested work deadlines and midterm dates.



\subsection{Inputs and Outputs}

\textbf{Inputs:} 
\begin{enumerate}
    \item \textbf{Course Outlines:} The system will take detailed course outlines from the current
    academic year as input, including assignment deadlines, midterm dates, and
    course learning objectives. These outlines will be used to identify scheduling
    conflicts, evaluate workload distribution across courses, and assess compliance
    with CEAB accreditation requirements. For testing, the system will also take
    past or mock course outlines as input. 
    \item \textbf{Academic Calendar:} The system will take the current and future academic calendar year as input to 
    take into consideration holidays, test and examination restrictions, and final examination periods, to ensure 
    midterms, assignments and exams are scheduled correctly according to school policy. 
\end{enumerate}

\textbf{Outputs:}
\begin{enumerate}
    % cspell:ignore CEAB
    \item \textbf{Shared Academic Calendar:} The system will provide a
    centralized, calendar-like interface displaying midterm dates and
    assignment deadlines across all CAS courses and years. Professors and administrators
    will be able to filter the calendar with specific
    courses and years to view relevant deadlines and identify potential scheduling
    conflicts.

    \item \textbf{CEAB Compliance Notifications:} The system will analyze
    course outlines against CEAB accreditation requirements and notify
    professors when their outlines do not meet the required criteria. These
    notifications will identify potential compliance issues and help professors
    address missing or incomplete information.

    \item \textbf{Scheduling Recommendation Feature:} The system will provide
    professors with suggestions for adjusting midterm and assignment dates to
    reduce scheduling conflicts and student workload during particularly busy
    weeks. Each recommendation will include an explanation of the reasoning
    behind the suggested changes, allowing professors to make informed
    scheduling decisions. Recommendations will be made instead of automated schedule 
    changes, to ensure human verification.
\end{enumerate}

\subsection{Stakeholders}

The key stakeholders for this project include:
\begin{itemize}
    \item \textbf{CAS Administration:} The primary stakeholders and end-users of
    the system. They will use it to manage and maintain course information, monitor
    assessment schedules across CAS courses, and streamline the review of course
    outlines for CEAB accreditation requirements. The CAS Administration staff (Matt Vonk \& Connie Carrabs) 
    will also be responsible for overseeing the system's development and ensuring that it meets
    administrative needs.

    \item \textbf{Faculty Professors and Course Instructors:} End-users who would use the system to 
    review assessment schedules across CAS courses and identify potential conflicts with other midterms and 
    assignments. They will also be able to receive scheduling suggestions that propose alternative 
    assessment dates and provide reasoning for how these changes could help reduce student workload 
    and assessment conflicts.

    \item \textbf{Academic Advising Staff:} End-users who would use the centralized assessment schedule 
    to better understand periods of high student workload and identify potential scheduling conflicts that 
    could affect students. This information can support academic advising and help staff provide students 
    with more informed guidance regarding their course schedules and academic workload.
\end{itemize}

\subsection{Environment}

\wss{Hardware and Software Environment for the users.  The developer environment
is summarized as part of the developer plan.}

\section{Goals}

The project goals include:
\begin{enumerate}
    \item \textbf{Centralized Assessment Calendar}

    Develop a shared calendar that allows professors and CAS administrators to view midterm and assignment deadlines across all supported CAS courses, with filters for course, year, and program. Success will be measured by successfully displaying at least 95\% of assessment dates from the course outlines in the test dataset and correctly filtering assessments by course, year, and program.

    \item \textbf{Automated Course Outline Data Extraction}

    Automatically extract assessment dates, weights, and learning outcomes from course outlines and store them in a centralized CAS database for administrators to review, correct, and export. The system will flag missing, unclear, or incomplete information for human verification. Success will be measured by achieving at least 90\% extraction accuracy for assessment dates and weights and correctly flagging at least 90\% of predefined missing or incomplete fields in the test dataset.

    \item \textbf{Assessment Workload and Conflict Detection}

    Identify weeks with high concentrations of assessments for students in the same year and program, as well as potential conflicts with exam periods and relevant scheduling policies. Success will be measured by correctly identifying at least 90\% of predefined scheduling conflicts in a test dataset and ensuring that all identified conflicts are associated with the correct student cohort.

    \item \textbf{Scheduling Recommendations}

    Provide professors with alternative assessment dates that aim to reduce workload clustering and scheduling conflicts. Each recommendation will include an explanation of its reasoning and expected impact, while professors retain full control over whether to accept the suggested changes. Success will be measured by generating at least two feasible alternative dates for 90\% of test cases with scheduling conflicts, with at least one alternative reducing the number of identified conflicts in each applicable case.

    \item \textbf{Automated Accreditation and Policy Compliance Checking}

    Evaluate course outlines against McMaster's course outline policies and applicable CEAB accreditation requirements, identifying missing information and potential compliance gaps. Notify professors and administrators of issues requiring attention. Success will be measured by correctly identifying at least 90\% of predefined compliance issues in a test dataset and providing an actionable notification for each detected issue.
\end{enumerate}

\section{Stretch Goals}

\wss{Goals you hope to get to, but not necessary.  They are also helpful if the
scope of the original project needs to be expanded.}

\section{Custom Documents (CDs)}

\begin{comment}
\wss{For CAS 741: State whether the project is a research project. This
designation, with the approval (or request) of the instructor, can be modified
over the course of the term.}

\wss{For SE Capstone: List your custom documents (CDs).  Potential CDs include
usability testing, code walkthroughs, user documentation, formal proof,
GenderMag personas, Design Thinking, etc.  (The full list is on the course
outline and in Lecture 02.) Normally the number of CDs will be two (one per
term).  Approval of the CDs will be part of the discussion with the instructor
for approving the project. The CDs, with the approval (or request) of the
instructor, can be modified over the course of the term.}
\end{comment}

The following two Custom Documents were selected to guide our design process
throughout the year and support end-user adoption:

\begin{enumerate}
    \item \textbf{Fall Term CD:} \href{https://github.com/smiths/capTemplate/tree/main/docs/CDs/Wireframe}{Wireframe Report (WIREFRAME)}
    

    The proposed system centers on a visual administrative dashboard featuring
    a shared assessment calendar, course filtering tools, and scheduling recommendation
    displays for instructors and CAS administrators. Developing user interface mockups during the
    Fall term will allow us to define and visualize the UI component interactions and workflow navigation
    before beginning front-end development. This will also ensure that we are aligning our design with the needs of the users, 
    as well as receiving early feedback on the intuitiveness and usability of the system interface from our clients and stakeholders.

    \item \textbf{Winter Term CD:} \href{https://github.com/smiths/capTemplate/tree/main/docs/CDs/UserManual}{User Manual (USERMAN)}
    
    
    To ensure a smooth and seamless adoption process by instructors and the CAS administrative staff, 
    a comprehensive user manual will be created to provide clear, step-by-step guidance on operating the software. 
    The manual will cover uploading course outlines, filtering the centralized assessment calendar, 
    exporting database records, and reviewing automated policy compliance flags and scheduling recommendations.
\end{enumerate}


\newpage{}

\section*{Appendix --- Reflection}

\wss{Not required for CAS 741}

\input{../Reflection.text}

\begin{enumerate}
    \item What went well while writing this deliverable? 
    \item What pain points did you experience during this deliverable, and how
    did you resolve them?
    \item How did you and your team adjust the scope of your goals to ensure
    they are suitable for a Capstone project (not overly ambitious but also of
    appropriate complexity for a senior design project)?
\end{enumerate}  

\end{document}